\section{Расчётная часть}

\subsection{Исходные данные (вариант \varNumber)}

Расчёт выполняется для подшипника с геометрией~\varGeomLetter{} и микрорельефом типа~\varDimpleType{} (\varDimpleTypeName{}).

\subsubsection{Геометрические параметры}

\begin{itemize}
    \item Радиус вала: $R = \varR$~м (\varRmm~мм)
    \item Радиальный зазор: $c = \varC$~м (\varCmkm~мкм)
    \item Длина подшипника: $L = \varL$~м (\varLmm~мм)
    \item $L/D = \varLD$
\end{itemize}

\subsubsection{Параметры микрорельефа (\varDimpleTypeName)}

Профиль углубления:
\begin{equation}
    \varDimpleFormula
\end{equation}

Параметры: $a = \varDimpleAmm$~мм, $b = \varDimpleBmm$~мм, $h_p = \varDimpleHpmkm$~мкм ($h_p/c = \varDimpleHpRatio$). Размещение: $N_\varphi = \varNphi$, $N_Z = \varNZ$, диапазон \varPhiRange.

\subsubsection{Режимы работы}

$n = \varN$~об/мин, $\eta = \varEta$~Па$\cdot$с, $\varepsilon = 0.05 \div 0.8$. Сетка $500 \times 500$, $\omega = 1.5$, $\delta < 10^{-5}$.

\subsection{Результаты для гладкого подшипника}

% TODO: раскомментировать \includegraphics после генерации графиков

\begin{figure}[H]
\centering
%\includegraphics[width=0.85\textwidth]{figures/smooth_pressure_2d.png}
\fbox{\parbox{0.8\textwidth}{\centering\vspace{2.5cm}Графики будут вставлены после расчёта\vspace{2.5cm}}}
\caption{Распределение давления $P(\varphi)$ при $Z=0$ для гладкого подшипника ($\varepsilon = 0.6$)}
\end{figure}

\begin{table}[H]
\centering
\caption{Характеристики гладкого подшипника}
\begin{tabular}{|c|c|c|c|}
\hline
$\varepsilon$ & $F$, Н & $\mu$ & $Q$, л/с \\
\hline
0.05 & --- & --- & --- \\
\hline
0.2 & --- & --- & --- \\
\hline
0.4 & --- & --- & --- \\
\hline
0.6 & --- & --- & --- \\
\hline
0.8 & --- & --- & --- \\
\hline
\end{tabular}
\end{table}

\subsection{Результаты для подшипника с микрорельефом}

\begin{figure}[H]
\centering
%\includegraphics[width=0.85\textwidth]{figures/textured_clearance_2d.png}
\fbox{\parbox{0.8\textwidth}{\centering\vspace{2.5cm}Графики будут вставлены после расчёта\vspace{2.5cm}}}
\caption{Распределение зазора $H(\varphi)$ при $Z=0$ (\varDimpleTypeName, $\varepsilon = 0.6$)}
\end{figure}

\begin{figure}[H]
\centering
%\includegraphics[width=0.85\textwidth]{figures/example_pressure_3d.png}
\fbox{\parbox{0.8\textwidth}{\centering\vspace{3cm}3D-график давления\vspace{3cm}}}
\caption{Трёхмерное распределение давления $P(\varphi, Z)$ ($\varepsilon = 0.6$)}
\end{figure}

\subsection{Сравнительный анализ}

\begin{table}[H]
\centering
\caption{Сравнение характеристик при $\varepsilon = 0.6$}
\begin{tabular}{|l|c|c|c|}
\hline
\textbf{Характеристика} & \textbf{Гладкий} & \textbf{С углублениями} & \textbf{$\Delta$, \%} \\
\hline
Нагрузка $F$, Н & --- & --- & --- \\
\hline
Коэфф. трения $\mu$ & --- & --- & --- \\
\hline
Расход $Q$, л/с & --- & --- & --- \\
\hline
\end{tabular}
\end{table}

\begin{figure}[H]
\centering
%\includegraphics[width=0.85\textwidth]{figures/example_load.png}
\fbox{\parbox{0.8\textwidth}{\centering\vspace{2.5cm}$F(\varepsilon)$\vspace{2.5cm}}}
\caption{Нагрузочная способность в зависимости от эксцентриситета}
\end{figure}

\begin{figure}[H]
\centering
%\includegraphics[width=0.85\textwidth]{figures/example_friction.png}
\fbox{\parbox{0.8\textwidth}{\centering\vspace{2.5cm}$\mu(\varepsilon)$\vspace{2.5cm}}}
\caption{Коэффициент трения в зависимости от эксцентриситета}
\end{figure}

\begin{figure}[H]
\centering
%\includegraphics[width=0.85\textwidth]{figures/example_flow.png}
\fbox{\parbox{0.8\textwidth}{\centering\vspace{2.5cm}$Q(\varepsilon)$\vspace{2.5cm}}}
\caption{Расход смазки в зависимости от эксцентриситета}
\end{figure}

\subsubsection{Физическая интерпретация}

Нанесение \varDimpleTypeGenitive{} углублений приводит к увеличению несущей способности (микрогидродинамический эффект), снижению коэффициента трения (увеличение средней толщины плёнки) и росту расхода смазки (улучшение циркуляции и теплоотвода).

\newpage
